\documentclass[12pt,a4paper]{article}

\usepackage[english,italian]{babel}
\usepackage[T1]{fontenc}

\usepackage[margin=2cm]{geometry} %margini
\usepackage{float} %float per immagini e tabelle
\usepackage{graphicx} %figure
\usepackage{minted}

\setcounter{secnumdepth}{4}

\newcommand{\htmlcode}[1]{\mintinline{html}{#1}}
\newcommand{\csscode}[1]{\mintinline{css}{#1}}
\newcommand{\jscode}[1]{\mintinline{js}{#1}}
\newcommand{\phpcode}[1]{\mintinline{php}{#1}}

\usepackage{hyperref}

\title{Progetto di Tecnologie Web}
\author{}
\date{Anno accademico 2019/2020}

\begin{document}
\pagenumbering{gobble}
\maketitle
\begin{figure}[H]
	\centering
	\includegraphics[width=8cm]{img/logo/coque_scritta.png}
\end{figure}
\begin{table}[H]
	\centering
	\begin{tabular}{c|c c c}
		\textbf{Componenti} & Agatea     & Riccardo  & 1170718 \\
		                    & Bosinceanu & Ecaterina & 1169669 \\
		                    & Righetto   & Sara      & 1174009 \\
		                    & Schiavon   & Rebecca   & 1163774 \\
	\end{tabular}
\end{table}

\begin{center}
	\textbf{Indirizzo sito web}: http://tecweb1920.studenti.math.unipd.it/ragatea\\
	\textbf{Email referente del gruppo}: riccardo.agatea@studenti.unipd.it
\end{center}

\begin{table}[H]
	\centering
	\begin{tabular}{c|c c}
		\textbf{Utenti} & \textbf{E-mail} & \textbf{Password} \\
		\hline
		Admin           & admin@gmail.com & admin             \\
		User            & user@gmail.com  & user              \\
	\end{tabular}
\end{table}
\newpage
\pagenumbering{roman}
\tableofcontents
\newpage
\pagenumbering{arabic}
\renewcommand{\abstractname}{Abstract}
\begin{abstract}
	Il sito \emph{Uovo alla coque} è un sito di ricette a tema uovo, ovvero tutte le ricette al suo interno hanno come ingrediente principale o necessario l'uovo. Il nome del sito rispecchia questa caratteristica perchè rimanda alla ricetta più famosa con cui è possibile cucinare un uovo. \newline
	Il sito permette a tutti gli utenti di visualizzare tutte le ricette del sito, divise per portata (primi piatti, secondi piatti e dolci), e cercare le ricette che più si desiderano. Inoltre un utente registrato può commentare e votare le ricette preferite, mentre l'utente amministratore può anche aggiungere, modificare ed eliminare le ricette.
\end{abstract}
\newpage
\section{Fase di analisi}
\subsection{Analisi delle caratteristiche dell'utenza}
\label{sub:analisi_delle_caratteristiche_dell_utenza}
\paragraph{Destinatari}
\label{par:destinatari}

% paragraph destinatari (end)
\paragraph{Ricerca}
\label{par:ricerca}
Essendo rivolto anche a persone con poca dimestichezza nella navigazione in Internet, il sito deve offrire strumenti di ricerca efficaci ed intuitivi. L'idea più semplice è di inserire una barra di ricerca, visibile in ogni pagina, che permetta di raggiungere una qualunque pagina del sito a partire da alcune parole chiave.
% paragraph ricerca (end)
\paragraph{Accessibilità}
\label{par:accessibilità}
Non c'è nessun motivo per cui persone con disabilità varie, anche visive, abbiano precluse le informazioni offerte dal sito. 
% paragraph accessibilità (end)
% subsection analisi_delle_caratteristiche_dell_utenza (end)
\newpage
\section{Fase di progettazione}
\label{sec:fase_di_progettazione}
\subsection{Struttura}
\label{sub:struttura}
\subsubsection{Header}
\label{ssub:header}
L'header contiene logo, nav, breadcrumb e due link per il login e il signup. Il logo, situato in alto a sinistra, non rimanda a nessuna pagina in quanto il link per la homepage è già contenuto nel menu sotto la voce Home. I due pulsanti per il login e il signup sono rispettivamente ``Accedi'' e ``Registrati'' e sono entrambi presenti solo nel caso in cui l'utente non abbia ancora effettuato l'accesso. Una volta effettuato l'accesso, al posto dei due link c'è il nickname dell'utente e il link ``esci'' che scollega l'utente e lo porta alla home.
% subsubsection header (end)
\paragraph{Nav}
\label{par:nav}
Il nav del sito contiene il menu e la barra di ricerca.
Il menu occupa la posizione centrale dell'header e contiene i link alle pagine principali del sito:
\begin{itemize}
	\item Home page;
	\item Pagina relativa ai primi piatti;
	\item Pagina relativa ai secondi piatti;
	\item Pagina relativa ai dolci;
\end{itemize}
La barra di ricerca è situata sotto il menu in posizione centrale ed è ben visibile per permettere agli utenti di cercare le ricette di cui hanno bisogno.
% paragraph nav (end)
\paragraph{Breadcrumb}
\label{par:breadcrumb}
Il breadcrumb è posizionato a sinistra sotto la barra di ricerca e serve a identificare la posizione dell'utente all'interno del sito. L'ultimo campo corrisponde alla pagina corrente e, per evitare link circolari, è un testo.
% paragraph breadcrumb (end)
\subsubsection{Content}
\label{ssub:content}
Lo scopo principale del sito è cercare e consultare le ricette a cui si è interessati.

\paragraph{Pagina Home}
È la prima pagina ad essere visualizzata quando un utente visita il sito, contiene una breve descrizione dei contenuti.
Sono poi presenti le tre ricette (una per portata) con il voto più alto.

\paragraph{Pagine delle portate}
Le pagine con gli elenchi dei primi piatti, i secondi piatti e i dolci hanno la stessa struttura: in ogni pagina viene presentata la lista di tutte le ricette relative alla portata selezionata. Le ricette sono presentate in riquadri contenenti l'immagine del piatto, il nome, la difficoltà, il tempo necessario allo svolgimento, il voto medio e il link ``apri'' che porta alla pagina della ricetta selezionata.

\paragraph{Pagine delle ricette}
È la pagina in cui viene presentata la ricetta selezionata. Per prima cosa, è visibile il titolo della ricetta, l'immagine del piatto e sulla destra la lista degli ingredienti, la difficoltà e la durata della ricetta; in seguito è descritto il procedimento, seguito da un form che permette di assegnare alla ricetta un voto da 1 a 5. Infine c'è la sezione dei commenti in cui è possibile inserire un commento relativo alla ricetta e leggere i commenti degli altri utenti.

\paragraph{Pagine di Login e Signup}
L'utente ha la possibilità di registrarsi o accedere con le proprie credenziali tramite le pagine di login e signup. Nella pagina signup sono presenti dei campi per inserire l'indirizzo email, il nickname e la password (quest'ultima da inserire due volte per confermare la correttezza). Per il login è necessario inserire  il nickname e la password.

\paragraph{Pagina Utente}
Una volta effettuato l'accesso, l'utente può visitare la propria area personale e verificare le proprie informazioni correnti. C'è la possibilità di eliminare l'account inserendo la propria password nel form apposito in fondo alla pagina.

\paragraph{Pagina Modifica Utente}
Dalla pagina utente è possibile aprire la pagina modifica utente in cui sono presenti tutti i form necessari per cambiare le informazioni dell'utente (immagine, nickname, email e password). Ciascun form permette la modifica di un unica informazione, e deve essere confermato inserendo la password attuale. Nel caso si vogliano annullare le modifiche, lo si può fare mediante il pulsante ``Annulla'', in fondo alla pagina.

\paragraph{Pagine Aggiungi e Modifica ricetta}
Queste pagine sono raggiungibili solo dagli utenti con permessi amministratore. In entrambe le pagine sono presenti i form per l'inserimento delle informazioni che poi compariranno nella pagina della ricetta. Nella pagina ``modifica ricetta'' i form sono già compilati con le informazioni attuali. Anche in questo caso, per annullare le modifiche è presente il pulsante ``Annulla'' in fondo alla pagina.

% subsubsection content (end)
\subsubsection{Footer}
\label{ssub:footer}
Il footer contiene la dicitura per il copyright, i nomi degli autori del sito, e le icone di certificazione della validità rispetto agli standard W3C di XHTML e CSS3.

% subsubsection footer (end)
% subsection struttura (end)
\subsection{Visualizzazione} %le pagine devono essere accessibili indipendentemente dalle dimensioni del dispositivo e del browser
\label{sub:visualizzazione}

\subsubsection{Desktop}
\label{ssub:desktop}

% subsubsection desktop (end)
\subsubsection{Mobile}
\label{ssub:mobile}
Nella visualizzazione da mobile è stato implementato il menù ad hamburger, ovvero il menù è nascosto e per farlo comparire serve premere il burger menù situato in alto a destra. Il pulsante è facilmente raggiungibile per gli utenti che usano il telefono con due mani, mentre gli utenti destri che usano il telefano con una sola mano potranno aprire comodamente il menù col pollice.
Una volta aperto, il menù occupa quasi tutto lo schermo con le voci al centro. Sono presenti inoltre i link per il login e il signup e la barra di ricerca posizionata sul fondo del menù. Si è deciso di inserire tutte le voci insieme alla barra di ricerca nel menù, per avere tutti gli elementi necessari per la ricerca raccolti in un unico spazio. Quando il menù è chiuso è visibile solo il logo affiancato dal tasto del menù.


% subsubsection mobile (end)
\subsubsection{Stampa}
\label{ssub:stampa}
Per quanto riguarda la stampa, sono stati eliminati gli elementi di presentazione (immagini e background) e sono stati tenuti solo gli elemeti necessari, ovvero il contenuto della pagina. È stato rimosso il menù, la barra di ricerca e il footer in quanto non necessari per la visualizzazione della pagina stampata, mentre sono stati tenuti il breadcrumb per specificare il percorso alla pagina selezionata, e il logo per distinguere il sito.
Le pagine sono in bianco e nero per dare valore al contenuto anzichè alla presentazione.
% subsubsection stampa (end)
% subsection visualizzazione (end)

\subsection{Accessibilità}
\label{sub:accessibilità}
% come da specifiche è stata rispettata la completa separazione tra comportamento, presentazione e struttura.
% ricerca: già scritta nell analisi
% torna su: poichè il menu sia nella visualizzazione desktop che quella mobile, è fisso all'inizio di ogni pagina, per raggiungerlo facilmente dopo uno scroll verticale, sono stati posizionati delle ancore "torna su", che rimandano all' header della pagina attuale.
% il logo nell'header non rimanda alla home
% Dare le informazioni con un numero limitato di click: con solo due click si può accedere ad una ricetta, per esempio Primi piatti > Carbonara. 
% Non obbligare utente e reinserire i campi dei form dopo messaggi di errore. -> session di php
% Ogni messaggio di errore deve essere corredato da "che errore è" e da "Come fare a risolvere". validazione php js
% tutte le immagini hanno specificato la proprietà alt
% form
% non c'è testo libero, tutto è incapsulato dentro ad appositi tag descrittivi

\subsubsection{Convenzioni interne}
\label{ssub:convenzioni_interne}
Una delle convenzioni interne al sito riguarda i link: sono sottolineati per permettere l'immediata individuazione, e quelli visitati diventano gialli. % non ci sono link circolari. Quando sei loggato, il registrati nel contenuto della home non è più cliccabile
% subsubsection convenzioni_interne (end)
% fare riferimento al w3c https://www.w3.org/standards/webdesign/accessibility
% subsection accessibilità (end)

% section fase_di_progettazione (end)

\newpage
\section{Fase di implementazione}
\label{sec:fase_di_implementazione}
\subsection{Linguaggi}
\subsubsection{HTML}
Tutte le pagine del sito sono state realizzate utilizzando come linguaggio di markup XHTML piuttosto che HTML5, per garantire un certo grado di qualità del codice e alta compatibilità con i browser più obsoleti. In particolare le differenze più importanti che ci hanno portato a preferire XHTML ad HTML5 sono:
\begin{itemize}
	\item I tag \htmlcode{<html>}, \htmlcode{<head>}, \htmlcode{<title>} e \htmlcode{<body>} sono obbligatori;
	\item Gli elementi devono essere nidificati correttamente;
	\item Gli elementi devono essere sempre chiusi;
	\item Gli elementi devono essere sempre in minuscolo;
	\item I nomi degli attributi devono essere sempre in minuscolo;
\end{itemize}
Per le form di registrazione e accesso, abbiamo scelto di utilizzare come input per le e-mail il tipo \htmlcode{text}, che grazie ad opportuni controlli specificati nelle sezioni \S\ref{subs:php} e \S\ref{subs:javascript}, si comporta quasi similmente al tipo \htmlcode{email} di HTML5. Allo stesso modo, anche gli input che riguardano il tempo di una ricetta, la difficoltà di una ricetta e la barra di ricerca, sono tutti di tipo \htmlcode{text} al posto dei tipi \htmlcode{number}, \htmlcode{range} e \htmlcode{search}.

\subsubsection{CSS}

\subsubsection{SQL}
Il database mantiene le informazioni strutturali sulle ricette, comprese metainformazioni come l'autore della pagina, e le informazioni di accesso sugli utenti.
Inoltre, mantiene i voti e i commenti degli utenti.
Le chiavi primarie delle tabelle sono degli identificatori numerici autoincrementali, eccetto per \sqlcode{voti}, per la quale la chiave primaria è data dall'utente che ha dato il voto e la ricetta che è stata votata, entrambi chiavi esterne.
Essendo plausibili più commenti dello stesso utente alla stessa ricetta, non abbiamo potuto utilizzare la stessa soluzione per la tabella \sqlcode{commenti}, per la quale sarebbe stato necessario includere anche l'orario, possibilità preclusa da \textit{MariaDB}.

Oltre alle chiavi primarie, abbiamo imposto un vincolo di unicità sul nome delle ricette e su nome utente ed email degli utenti, per evitare inutili confusioni in caso di ripetizioni.
Per prevenire riferimenti a contenuti rimossi, all'eliminazione di una ricetta o un utente vengono eliminati anche i rispettivi commenti e voti.

\subsubsection{PHP}\label{subs:php}
\paragraph{Costruzione delle pagine}
Tutte le pagine del sito sono costruite dinamicamente.
Abbiamo definito la struttura essenziale comune a tutte le pagine attraverso un template, utilizzando degli elementi vuoti invalidi (i cui tag sono riconoscibili immediatamente grazie al suffisso \htmlcode{Placeholder}) per rappresentare porzioni dell'albero XML diverse fra le pagine.
Queste porzioni dinamiche sono definite in appositi file separati; le funzioni \phpcode{file_get_contents($filename)} e \phpcode{str_replace($search,$replace,$subject)} permettono di inserire il contenuto di tali file in sostituzione ai tag invalidi, ottenendo un documento completo e valido.
In alcuni casi, per semplificare la costruzione delle pagine, anche i file parziali contengono tag invalidi, i quali vengono sostituiti applicando ricorsivamente la stessa procedura.
Abbiamo scelto di utilizzare l'estensione \textit{.php} per tutti i file contenenti tag invalidi, per evidenziare la loro non completezza e non validità.
Fra i tag invalidi, due sono di particolare rilevanza:
\begin{itemize}
	\item \htmlcode{<rootFolder />} viene utilizzato per indicare un path relativo alla pagina corrente della radice del sito.
	\item \htmlcode{<backToTopPlaceholder />} indica la posizione di un link \textit{Torna Su}.
\end{itemize}
Questi due tag possono essere utilizzati in tutti i file, e vengono sostituiti con il rispettivo codice valido solo al termine di tutte le altre sostituzioni.

Per una comoda gestione del template principale, abbiamo definito una classe \phpcode{DBConnection}.
Ogni script che deve ritornare una pagina al client istanzia questa classe, ed invocando i suoi metodi la costruisce.
Sono disponibili i metodi:
\begin{itemize}
	\item \phpcode{__construct($root,$standard)} per costruire un'istanza della classe. \phpcode{$root} deve essere il path relativo alla directory corrente della directory radice del sito, cioè quella contenente il file \texttt{index.php}. \phpcode{$standard} deve essere la stringa \phpcode{"xhtml"}. In origine era prevista la possibilità di scelta fra gli standard XHTML e HTML5 in base alla pagina attraverso questo parametro; dopo avere scelto di utilizzare solo XHTML, abbiamo comunque mantenuto questo parametro a fini di  manutenibilità ed estensibilità.
	\item \phpcode{setTitle($title)} per inserire il titolo della pagina, sia nel tag \htmlcode{title} sia nel tag \htmlcode{meta} con \htmlcode{name="title"}.
	\item \phpcode{setDescription($description)} per inserire la descrizione della pagina nel tag \htmlcode{meta} con \htmlcode{name="description"}.
	\item \phpcode{setAuthor($author)} per inserire il nome dell'autore della pagina nel tag \htmlcode{meta} con \htmlcode{name="author"}.
	\item \phpcode{setOtherMeta($tags)} per inserire altri tag \htmlcode{meta}.
	\item \phpcode{setLogin($loginSnippet)} per inserire i link alle pagine di login e signup se l'utente non è autenticato, oppure un messaggio di benvenuto, il link alla pagina utente ed il link per il logout se è autenticato.
	\item \phpcode{setNav($navSnippet)} per inserire il menu, i cui link possono venire rimossi per evitare riferimenti circolari.
	\item \phpcode{setBreadcrumb($breadcrumbSnippet)} per inserire il breadcrumb specifico per ogni pagina.
	\item \phpcode{setContent($contentSnippet)} per inserire il contenuto della pagina.
	\item \phpcode{setBackToTop($backToTopSnippet)} per inserire i link \textit{Torna Su}.
	\item \phpcode{setAnnulla($annullaSnippet)} per inserire i pulsanti \textit{Annulla}.
	\item \phpcode{setFooter($footerSnippet)} per inserire il footer del sito
	\item \phpcode{send()} per completare l'elaborazione ed inviare la pagina completa al client.
	\item \phpcode{save($file)} per completare l'elaborazione e salvare la pagina completa su file.
\end{itemize}

\paragraph{Pagine di errore}

Per segnalare all'utente errori nelle richieste al server, causati ad esempio dall'inserimento di un URL incorretto, abbiamo implementato delle pagine di errore corrispondenti ai codici HTTP rilevanti per il sito, in particolare per gli errori lato client.
Ogni pagina di errore è implementata in un file nella radice del sito, nominato \texttt{XXX.php}, dove \texttt{XXX} è il rispettivo codice HTTP\@.

\begin{itemize}
	\item Errore \textit{400 BAD REQUEST}: ritornato nel caso l'URL della richiesta non rispetti vincoli di coerenza, ad esempio l'esecuzione di una ricerca priva del termine della ricerca nella query.
	\item Errore \textit{401 UNAUTHORIZED}: ritornato nel caso in cui sia necessaria un'autenticazione non fornita, ad esempio cercando di inserire un commento prima di aver eseguito il login.
	\item Errore \textit{403 FORBIDDEN}: ritornato nel caso in cui l'utente sia autenticato, ma non abbia i permessi necessari ad eseguire l'operazione richiesta, ad esempio cercando di inserire una nuova ricetta senza essere un amministratore.
	\item Errore \textit{404 NOT FOUND}: ritornato nel caso in cui la richiesta abbia una sintassi non riconosciuta, ad esempio se l'URL individua un file non esistente.
\end{itemize}

% funzionalità di ricerca
% query portata
\paragraph{Validazione lato server}\label{par:validazione_lato_server}
Le funzioni di validazione degli input lato server sono contenute all'interno del file \texttt{validation.php}. I controlli implementati si comportano allo stesso modo di quelli descritti nella sezione \S\ref{subs:javascript} aggiungendo loro il controllo anche dei dati presenti nel database, nello specifico:
\begin{itemize}
	\item \phpcode{checkNickname($stringNickname)} che oltre a controllare se l'input inserito dall'utente contiene non più di 20 caratteri alfanumerici, controlla anche che il nome non sia già presente nel database.
	\item \phpcode{checkEmail($stringEmail)} che valida la mail grazie al filtro offerto da PHP e inoltre controlla che la mail non sia già presente nel database.
	\item \phpcode{comparePassword($stringPassword, $stringPasswordConfirm)} che verifica, se entrambi i campi non sono vuoti, che le due password siano uguali tra loro.
	\item \phpcode{checkLogin($stringPassword, $stringNickname)} controlla che sia il nickname che la password inseriti dall'utente non siano stringhe vuote, che il nickname inserito sia presente nel database e che ad esso corrisponda la passowrd immessa dall'utente.
	\item \phpcode{checkNomeRicetta($stringNomeRicetta, $dbCondition)} che varifica se la stringa rispetta la lunghezza di massimo 55 caratteri alfanumerici e, se la variabile \phpcode{dbCondition} è true, anche che il nome non sia già presente nel database. Quest'ultimo controllo viene fatto solo nel caso si sia aggiungendo una nuova ricetta, mentre non viene fatto se si vuole modificare una ricetta senza cambiarne il nome.
	\item \phpcode{checkDifficolta($stringDifficolta)} controlla che la stringa che rappresenta la difficoltà sia un numero tra 1 e 5.
	\item \phpcode{checkTempo($stringTempo)} verifica che il tempo non sia vuoto e che sia un intero positivo, escluso lo zero.
	\item \phpcode{checkKeywords($stringKeywords)} controlla che la stringa inserita dall'utente non sia vuota e che contenga almeno un carattere alfanumerico.
	\item \phpcode{checkImage($stringImage)} verifica che il file caricato sia effettivamente un'immagine controllando la sua estensione, e che la sua dimensione non superi 150KB.
	\item \phpcode{checkCommento($testo)} che controlla grazie all'espressione regolare, che il testo del commento non sia vuoto. Non è stato utilizzato la funzione \phpcode{empty} perchè la stringa \phpcode{"0"} veniva considerata come vuota, mentre contiene un carattere.
\end{itemize}
Inoltre è stata implementata una funzione ausiliaria \phpcode{test_input($data)} che rimuove sia gli spazi che gli slash di escape, e trasforma in html i caratteri speciali, per esempio \phpcode{à} diventa \phpcode{&agrave;}.

\subsubsection{Javascript}\label{subs:javascript}
Il linguaggio Javascript è stato utilizzato per implementare il comportamento dinamico lato client di alcune pagine del sito, nello specifico il menu ad hamburger presente nella visualizzazione mobile e la validazione lato client delle form presenti nelle pagine di accesso, registrazione, modifica dati utente, aggiunta e modifica di una ricetta, aggiunta e modifica di un commento.
Poiché è possibile che non tutti gli utenti dispongano della tecnologia adatta, le funzionalità implementate in Javascript sono minime.
\paragraph{Menu mobile}\label{par:menu_mobile}
\paragraph{Validazione lato client}\label{par:validazione_lato_client}
Per quanto riguarda la validazione della form relativa alla registrazione, sono stati implementati i seguenti controlli:
\begin{itemize}
	\item \jscode{checkNickname(nickname)} che, passato come parametro la stringa che rappresenta l'id relativo al nickname presente nella form in esame, controlla se l'input inserito dall'utente è una stringa vuota e se non corrisponde all'espressione regolare che identifica stringhe contenti solo lettere e numeri di lunghezza tra 3 e 20 caratteri.
	\item \jscode{checkEmail(email)} che, passato come parametro la stringa che rappresenta l'id relativo alla mail presente nella form in esame, verifica sia se l'input inserito dall'utente è una stringa vuota, sia se non corrisponde all'espressione regolare che identifica le e-mail.
	\item \jscode{isPasswordEqual(password, passwordConfirm)} che, passati come parametri le stringhe gli id che identificano l'input relativo alla password e la conferma della password, controlla se le due password sono uguali tra loro.
\end{itemize}
Per la validazione degli input nella form di login, si è deciso di controllare grazie alla funzione \jscode{loginValidator(campo)} che entrambi i campi passati come paramentri, ossia il nickname e la password inseriti, non siano vuoti. \newline
Per la validazione delle form presenti nella pagina di modifica dati dell'utente, sono stati riutilizzati i controlli \jscode{checkNickname(nickname)}, \jscode{checkEmail(email)}, \texttt{isPasswordEqual(password, passwordConfirm)} ed è stato implementato il controllo per la validazione dell'immagine:
\begin{itemize}
	\item \jscode{checkImage(image)} che, data la stringa che rappresenta l'id per l'inserimento del file, verifica se l'estensione del file appartiene alle estensioni delle immagini accettate e se la dimensione del file è minore di 150KB.
\end{itemize}
Per le form delle pagine di aggiunta e modifica di una ricetta è stato riutilizzato il controllo sull'immagine e inoltre sono stati implementati i seguenti controlli:
\begin{itemize}
	\item \jscode{checkTitolo(titolo)} che, passato come parametro la stringa che rappresenta l'id relativo al titolo di una ricetta, verifica se l'input inserito dall'utente è una stringa vuota e se la lunghezza della stringa è tra 3 e 55 caratteri.
	\item \jscode{checkDifficolta(difficolta)} che, passato come parametro la stringa che rappresenta l'id relativo alla difficoltà di una ricetta, controlla se l'input inserito dall'utente è una stringa vuota e se corrisponde ad un numero tra 1 e 5.
	\item \jscode{checkTempo(tempo)} che, passato come parametro la stringa che rappresenta l'id relativo al titolo di una ricetta, controlla se l'input inserito dall'utente è una stringa vuota e se corrisponde ad un intero positivo, escluso lo 0.
	\item \jscode{checkKeywords(keywords)} che, data la stringa che identifica l'input del commento, controlla se contiene solo caratteri alfanumerici e che non sia vuoto.
\end{itemize}
Per quanto riguarda la validazione nell'inserimento e modifica dei commenti, è stata implementata la funzione \jscode{checkCommento(commento)} che, passato come parametro la stringa che identifica l'input della form, verifica che il commento non sia vuoto. \newline
Per la stampa degli eventuali messaggi di errore che possono incorrere in seguito alla validazione dei campi, è stato creato un metodo generico \jscode{printError(condition, id, message)} che passato come parametri \jscode{condition}, ossia se il campo soddisfa i controlli effettuati, \jscode{id} che identifica l'input del campo controllato, di cui verrà aggiunto in coda "-message" che identifica lo span dove stampare il messaggio d'errore e \jscode{message} che contiene il messaggio da stampare.
Nel caso in cui Javascript non sia disponibile o sia disabilitato, la validazione degli input è effettuata lato server grazie alle funzioni PHP sopra descritte.

% section fase_di_implementazione (end)

\newpage
\section{Fase di testing}
\label{sec:fase_di_testing}
\subsection{Validatori}
\label{sub:validatori}

% subsection validatori (end)
% section fase_di_testing (end)
\newpage
\section{Organizzazione interna}
% ripartizione dei compiti, organizzazione e ore
Affinchè ogni membro apprenda ogni aspetto e argomento nella realizzazione di un sito, si è deciso che ognuno debba implementare parti diverse del sito. Questo ha portato allo scambio di idee, opinioni e alla realizzazione di decisioni più pensate e consapevoli. Per contro, il gruppo ha speso più tempo nella propria formazione. \newline
Nello specifico l'analisi e la progettazione del sito è stata pensata in gruppo, mentre l'implementazione del sito è stata suddivisa tra i membri in questo modo:
\begin{itemize}
	\item Agatea Riccardo ha realizzato alcune pagine PHP nello specifico il template, la pagina di home e la funzionalità di ricerca.
	\item Bosinceanu Ecaterina ha realizzato alcune pagine PHP nello specifico la pagina di ricetta, la pagina di utente e modifica utente, le funzionalità di validazione delle ricette. % css desktop, js validazione delle ricette, stesura della relazione
	\item Righetto Sara ha realizzato alcune pagine PHP nello specifico le pagine di registrazione e accedi, la pagina per aggiungere una ricetta, la pagina per la modifica di una ricetta, la funzionalità di eliminazione di una ricetta e le funzionalità di validazione dell'accesso, registrazione e utente. % css mobile un pochettino, js validazione del resto, stesura della relazione
	\item Schiavon Rebecca ha realizzato alcune pagine PHP nello specifico la pagina generica degli elenchi che comprende \textit{primi piatti}, \textit{secondi piatti} e \textit{dolci}, le funzionalità relativa ai piatti più votati e la paginazione degli elenchi. % css un po tutti
\end{itemize}

% all'inizio è stato fatto il template, poi le pagine php/html aggiungendo il css ed infine il js
\end{document}
